\section*{Constant Degree Condensers (original Section 3)}
In this section we prove:
\begin{theorem}
For every constant $0<\delta_1<\delta_2=0.7$, every $s\ge s_0(\delta_1)$ and every integer $n_1$ and $\varepsilon_{\rm cond}\ge2^{-\Omega(n_1)}$ there exists an explicit $\delta_1\to_{\varepsilon_{\rm cond}}\delta_2$ condenser $C:\{0,1\}^{n_1}\times\{0,1\}^d\to\{0,1\}^{n_2}$ with $n_2=(\frac23)^sn_1$ and $d=4c_{\rm DKSS}(s+\log\frac1{\varepsilon_{\rm cond}})$, where $c_{\rm DKSS}$ is the constant from Theorem 15. Note that $d$ is independent of $n_1$.
\end{theorem}
Note that, in particular, for every $\delta_1>0$ there exists an explicit $\delta_1\to_{\varepsilon_{\rm cond}}\delta_2=0.7$ condenser $C:\{0,1\}^{n_1}\times\{0,1\}^d\to\{0,1\}^{n_2}$ with $n_2=\Omega(n_1)$ and $d=O(\log\frac1{\varepsilon_{\rm cond}})$. However, we will need the more precise version that appears in Theorem 23.
The proof goes through somewhere-random condensers, so let us first discuss the similarities and differences between condensers and somewhere-random condensers. We begin with the necessary definitions:
\begin{definition}[s.r. condenser]
A function $C:\{0,1\}^n\to(\{0,1\}^m)^A$ is an $(n,k)\to_\varepsilon(m,k')$ s.r. condenser if for every $(n,k)$ source $X$ it holds that $C(X)=C(X,1)\circ\ldots\circ C(X,A)$ is $\varepsilon$-close to a $k'$ s.r. source. If $k=\delta n$ and $k'=\delta m$ we say $C$ is $\delta\to_\varepsilon\delta'$ s.r. condenser.
\end{definition}
We may take a condenser $C:\{0,1\}^n\times\{0,1\}^d\to\{0,1\}^m$ and expand it to a table with the outputs of all possible seeds, i.e., define $S:\{0,1\}^n\to(\{0,1\}^m)^D$, with $D=2^d$, where $S(x)_i=C(x,i)$. The condenser property guarantees that for every $k$-source $X$, most rows in the table are close to having $k'$ min-entropy. In contrast, a s.r. condenser is a weaker object, because it only guarantees that one row has $k'$ entropy (or more precisely that we are in a convex combination of such cases).
The major question we consider now is the dependence of the degree ($2^d$ for condensers and $A$ for s.r. condensers) on $n,m,k,k'$ and $\varepsilon$. We focus on the case where $m=\Omega(n)$, $k=\delta n$, $k'=\delta'm$ and $\delta<\delta'$ are constants. A priori, we could have expected the degree to depend on $n$ and $\varepsilon$, as is indeed the case when $m$ might be arbitrarily small. However, remarkably, things are drastically different when $m=\Omega(n)$. In this case both condensers and s.r. condensers may be of degree that is independent of $n$ and this will be crucial for us. If we consider the dependence on the error, then s.r. condensers may have exponentially-small error and constant $D$, whereas the degree of a condenser is at least $d\ge\log(\frac1\varepsilon)$. Remarkably, all of that can be explicitly achieved, as we now explain.
The basic building block we use is the following beautiful result of Zuckerman, which is based on additive combinatorics:
\begin{theorem}[{[45, Theorem 8.3]}]
For every constant $0<c<1$ there exists a constant $\alpha=\alpha(c)$ such that for every constant $\delta\le c$ and integer $n$ there exists an explicit function $C:\{0,1\}^n\to(\{0,1\}^{\frac23n})^2$ that is a $\delta\to_\varepsilon(1+\alpha)\delta$ s.r. condenser with $\varepsilon=2^{-\Omega(\alpha\delta n)}$.
\end{theorem}
Somewhere-random condensers can be easily composed. Specifically, Barak et al. [5] showed that if $C_1:\{0,1\}^{n_1}\to(\{0,1\}^{n_2})^{\ell_1}$ is a $\delta_1\to_{\varepsilon_1}\delta_2$ s.r. condenser and $C_2:\{0,1\}^{n_2}\to(\{0,1\}^{n_3})^{\ell_2}$ a $\delta_2\to_{\varepsilon_2}\delta_3$ s.r. condenser then $C_2\circ C_1:\{0,1\}^{n_1}\to(\{0,1\}^{n_3})^{\ell_1\cdot\ell_2}$ defined by $C_2\circ C_1(x)_{(i_1,i_2)}=C_2(C_1(x)_{i_1})_{i_2}$ is a $\delta_1\to_{\varepsilon_1+\varepsilon_2}\delta_3$ s.r. condenser.
Composing the s.r condenser of Theorem 25 with itself $s$ times we get an explicit function $C:\{0,1\}^n\to(\{0,1\}^m)^D$ with $D=2^s$ and $m=(\frac23)^sn$ that is a $\delta\to_\varepsilon\delta'$ s.r. condenser with $\varepsilon=\sum_{i=1}^s2^{-\Omega((1+\alpha)^i\delta(\frac23)^in)}=2^{-\Omega(m)}$ and $\delta'\ge(1+\alpha(\delta'))^s\delta$. Therefore:
\begin{lemma}
For every constants $0<\delta_1<\delta_2<1$ there exists a constant $s=s(\delta_1,\delta_2)$ and an explicit function $C:\{0,1\}^{n_1}\to(\{0,1\}^{n_2})^D$ that is a $\delta_1\to_\varepsilon\delta_2$ s.r. condenser with $D=2^s$, $n_2=(\frac23)^sn_1$ and $\varepsilon=2^{-\Omega(n_2)}$. Note that $D$ is independent of $n$ and $\varepsilon$.
\end{lemma}
Right now, if $X$ is a $k$-source, the table $C(X)$ has $D$ rows, and, roughly speaking, the guarantee is that one of these rows has density $\delta'$. We want to change this to get a condenser, i.e., we are willing to invest a short seed (that is independent of $n$) and we want to get one output which is close to uniform. (Alternatively, we can write the condenser as a table with one row per seed, the number of rows is independent of $n$ and most rows are close to uniform.) This is exactly what a merger does and applying the merger of Theorem 15 with $\beta=\frac14$ on the s.r. condenser of Lemma 26 (with $\delta_2$ close to 1) gives Theorem 23.
\section*{The Unbalanced Two-Source Extractor Construction (original Section 4)}
The main result of this section is the following two-source extractor.
\begin{theorem}
For every integer $n$ and two constants $\delta_0,\varepsilon_0>0$ there exists a constant $c$ such that for $d\ge c\log n$ and $k\ge c\log d$ there exists an explicit $((n,k),(d,\delta_0d),\varepsilon_0)$ two-source extractor $\twoExt:\{0,1\}^n\times\{0,1\}^d\to\{0,1\}$.
\end{theorem}
The extractor in the above theorem has constant error, and works when:
\begin{enumerate}
\item Each source's entropy is in the order of the logarithm of the length of the other source.
\item The shorter source, of length $d$, has an arbitrarily small constant density $\delta_0$.
\end{enumerate}
We think of $n$ and $d=d(n)$ as growing parameters while $\varepsilon_0$ and $\delta_0$ are constants. We use asymptotic notations (such as $\Omega(\cdot)$) to hide constants that are independent of $n$ and $d$ (but may depend on $\varepsilon_0$ and $\delta_0$).
\subsection*{4.1 The Construction}
Recall that $\varepsilon_0$ is the target error of the extractor $\twoExt$. The input to $\twoExt$ is a pair $(x,y)$ where $x$ is drawn from an $(n,k)$ source $X$, and $y$ is drawn from an independent $(d,\delta_0d)$ source $Y$. Our problem is that the $y$ comes from a $\delta_0d$-source for some $\delta_0<\frac12$. To overcome this, we do the following:
\begin{itemize}
\item We apply the condenser of Theorem 23 on $y$ to get a table $y'$ that is 1-wise 0.7-dense. Notice that the output of this step is a table rather than a single output.
\item We apply Raz's extractor (Theorem 11) on the table and the input $x$ from the other source to convert the table $y'$ to another table $y''$ that is 1-wise uniform.
\item We apply the $t$ correlation-breaker with advice of Theorem 17 on $y$, using the table $y''$ as the seed, to get a table $y'''$ that is $t$-wise uniform.
\item Finally, we apply the resilient function $f$ of Theorem 22 on the table $y'''$ to collapse the many rows of the table to a single, close to uniform, output.
\end{itemize}
Formally, these steps work as follows:
\paragraph{Condense the short source.} We are given $\delta_0<\frac12$. Set $\delta'=0.69$ and $\delta_2=0.7$.
By Theorem 23 there exists a constant $s_0=s_0(\delta_0)$ such that for every $s\ge s_0$ there exists an explicit
\[C:\{0,1\}^d\times\{0,1\}^a\to\{0,1\}^{d'}\]
that is a $\delta_0\to_{\varepsilon_{\rm cond}}\delta_2=0.7$ condenser with $a=4c_{\rm DKSS}(s+\log\frac1{\varepsilon_{\rm cond}})$ and $d'=(\frac23)^sd$.
We set
\[\gamma=\frac1{2^5c_{\rm DKSS}},\]
and this also fixes $c_\gamma$ as in Theorem 22. Notice that $\gamma$ and $c_\gamma$ are fixed constants independent of all other parameters in our system.
Now, choose $\varepsilon_{\rm cond}$ so that
\begin{equation}\left(\frac1{\varepsilon_{\rm cond}}\right)^{\log(3/2)}\ge\frac4{\delta_0}2^{12}c_\gamma c_{\rm DKSS}^4\log^4\frac1{\varepsilon_{\rm cond}},\tag{1}\label{eq:cond}\end{equation}
and also so that $\varepsilon_{\rm cond}\le\xi(\varepsilon_0,\delta_0)$, where
\begin{equation}\xi(\varepsilon_0,\delta_0)=\min\left\{2^{-s_0},\left(\frac{\varepsilon_0}8\right)^2,\left(\frac{\varepsilon_0}5\right)^{c_\gamma},\left(\frac{\varepsilon_0}{5c_\gamma}\right)^{1/\gamma}\right\}.\tag{2}\label{eq:xi}\end{equation}
Given $\varepsilon_{\rm cond}$, we set $s=\log\frac1{\varepsilon_{\rm cond}}\ge s_0$, giving $a=8c_{\rm DKSS}\log\frac1{\varepsilon_{\rm cond}}$. Note that the degree of the condenser, $A=2^a$, satisfies
\[\sqrt{\varepsilon_{\rm cond}}A=2^{-\frac12\log\frac1{\varepsilon_{\rm cond}}+a}=2^{-\frac a{2^4c_{\rm DKSS}}+a}=A^{1-2\gamma}.\]
Observe that $s\ge s_0$ and that $d'=\Omega(d)$. Also, notice that $(\delta_2-\delta')d'=d'/100\ge a=\log A$ for large enough $d$. Thus, $C$ is a $(d,\delta d)\to_{\varepsilon_{\rm cond}}(d',\log(A)+\delta'd')$ condenser.
Define an $A\times d'$ table $Y'$ where
\[Y'_i=C(Y,i)\in\{0,1\}^{d'}\quad\text{for }i=1,\ldots,A.\]
\paragraph{1-wise uniformity.} Let $c_1,c_2$ be the constants from Theorem 11 for $\delta_{\rm Raz}=0.6$.
Notice that $\delta_{\rm Raz}d'=\Omega(d')=\Omega(d)$. Therefore, for a constant $c$ large enough, $d\ge c\log n$ is large enough so that $\delta_{\rm Raz}d'\ge c_2\log n$. We can, in particular, choose $c$ such that in addition $c\ge c_1$. Recalling that $k\ge c\log d$, we have $k\ge c_1\log d'$. By Theorem 11, there exists an explicit function
\[\Raz:\{0,1\}^n\times\{0,1\}^{d'}\to\{0,1\}^{d''}\]
that is a strong $((n,k),(d',\delta_{\rm Raz}d'),\varepsilon_{\rm Raz}=2^{-\Omega(d'')})$ two-source extractor with $d''=\Omega(\min\{k,\delta_{\rm Raz}d'\})=\Omega(k)$.\footnote{Although $k\ge c\log d$ we can always assume w.l.o.g. that $k=c\log d$ and so $k\le\delta_{\rm Raz}d'=\Omega(d)$.} Define an $A\times d''$ table $Y''$ where
\[Y''_i=\Raz(X,Y'_i)\quad\text{for }i=1,\ldots,A.\]
\paragraph{$t$-wise uniformity.} Let $k_{\rm CBA}=\frac{\delta_0d}8$ and $\varepsilon_{\rm CBA}=\frac1d$. Set
\[t=\frac{\delta_0}4\left(\frac32\right)^s.\]
Notice that for a large enough constant $c$ we have $d''=\Omega(k)=\Omega(c\log d)\ge c_{\rm CBA}at\log\frac d{\varepsilon_{\rm CBA}}$, where the latter is the seed-length required by the correlation-breaker from Theorem 17. Also, $k_{\rm CBA}=\frac{\delta_0d}8\ge d''$ for large enough $d$, as $d''=\Omega(k)=\Omega(\log d)$. Hence, by Theorem 17 there exists an explicit function
\[\CBA:\{0,1\}^d\times\{0,1\}^{d''}\to\{0,1\}\]
that is a $(t,k_{\rm CBA},\varepsilon_{\rm CBA})$ correlation-breaker with advice.
Define an $A\times1$ table $Y'''$ where
\[Y'''_i=\CBA(Y,Y''_i,i)\quad\text{for }i=1,\ldots,A.\]
Keep in mind that the entropy in $Y$ suffices for CBA since $H_\infty(Y)=8k_{\rm CBA}$.
\paragraph{Collapse.} Take $f:\{0,1\}^A\to\{0,1\}$ to be the $(q=A^{1-2\gamma},t,\varepsilon_f=c_\gamma A^{-\gamma})$ resilient function of Theorem 22 and output $f(y'''_1,\ldots,y'''_A)$.
\subsection*{4.2 Two Subtleties}
As mentioned in the introduction, there are several delicate issues in the analysis:
\begin{enumerate}
\item Circular dependence: $Y''$ depends on both $X$ and $Y$, and is used as a seed in the application of the correlation-breaker with advice on $Y$.
\item We need $Y'''$ to be close to a perfect $t$-wise independent table, while the correlation-breaker with advice only guarantees that every $t$ good rows are close to uniform. To bridge the gap we need the error to be at least polynomially-small in the number of rows, but some of the steps incur a large constant error.
\end{enumerate}
To overcome the first issue we show a conditioning under which $Y''$ is still good, $Y$ is independent of $Y''$ and even after the conditioning the two sources have enough min-entropy.
To overcome the second issue we distinguish between large errors that depend on the number of rows $A$, and small errors that depend on the row length (see Section 1.3 in the introduction of the original publication). In particular, the errors are of three types:
\begin{itemize}
\item The probability $p_1$ that a value we condition upon is bad. This error is incurred by the condenser and is high (think of it as being a constant).
\item We show that when we condition on a good value, every $t$ good rows in $Y'''$ are $p_2$-close to uniform. We then claim that $Y'''$ as a table is $A^tp_2$-close to a table where the good rows are truly $t$-wise independent (where $A$ is the number of rows in the table $Y'''$). The error $p_2$ is incurred by Raz's extractor and by the correlation-breaker, and can be made very small if we deal with a source $X$ having enough min-entropy. We make $p_2$ small enough so that $A^tp_2$ is also small.
\item A third error $p_3$ is incurred by the resilient function $f$. This error is large, say, a constant, and we are fine with that.
\end{itemize}
Note that we cannot just accumulate all errors as $A^tp_1$ is way larger than 1. Instead, we argue that with a constant probability $1-p_1$, we get extremely close to perfect behavior, and then we get such a small error $p_2$ so that $A^tp_2$ is also small.
